% Slide template for the FreeBSD Developer Summit.
% Take it or leave it :-)
%
% It requires LaTeX and LaTeX Beamer [1] to compile.
% pdfLaTeX is recommended for compilation as it produces
% PDF file immediately.
%
% $ pdflatex some.latex
%
%
% It is also recommended to convert images to PDF
% by using ImageMagick ("convert") before including them.
%
% [1] http://latex-beamer.sourceforge.net
%

\documentclass{beamer}

\usepackage{url}
\usepackage[utf8]{inputenc}
%\usepackage[T1,T2A]{fontenc}
\usepackage[T1]{fontenc}
\usepackage{textcomp}
%\usepackage[english,russian]{babel}
\usepackage[english]{babel}
\usepackage{verbatim}
\usepackage{graphicx}
\usepackage{listings}
\usepackage{mathtools}
\usepackage{color}
\usepackage{listings}
\usepackage{tikz}

\mode<presentation>
{
  \definecolor{beamer@gker}{rgb}{0.8,0.0,0.0}
  \setbeamercolor*{structure}{fg=beamer@gker}
  \logo{\includegraphics[scale=0.5]{../ganditemplate/logo.pdf}}
}

\setbeamertemplate{footline}[text line]{%
  \parbox{\linewidth}{\vspace*{-8pt}
  \prestitle
  \hfill
  \insertshorttitle
  \hfill
  \insertframenumber\ of \inserttotalframenumber}}
\setbeamertemplate{navigation symbols}{}

\definecolor{light-gray}{gray}{0.60}

\lstdefinestyle{custommake}{
  belowcaptionskip=1\baselineskip,
  breaklines=true,
  frame=L,
  xleftmargin=\parindent,
  language=Make,
  showstringspaces=false,
  basicstyle=\footnotesize\ttfamily,
  keywordstyle=\bfseries\color{green!40!black},
  commentstyle=\itshape\color{purple!40!black},
  identifierstyle=\color{blue},
  stringstyle=\color{orange},
}

\definecolor{mygreen}{rgb}{0,0.6,0}
\definecolor{mygray}{rgb}{0.5,0.5,0.5}
\definecolor{mymauve}{rgb}{0.58,0,0.82}



\lstset{ %
  backgroundcolor=\color{white},   % choose the background color; you must add \usepackage{color} or \usepackage{xcolor}
  basicstyle=\ttfamily\tiny,        % the size of the fonts that are used for the code
  breakatwhitespace=false,         % sets if automatic breaks should only happen at whitespace
  breaklines=true,                 % sets automatic line breaking
  captionpos=b,                    % sets the caption-position to bottom
  commentstyle=\color{mygreen},    % comment style
  deletekeywords={...},            % if you want to delete keywords from the given language
  escapeinside={\%*}{*)},          % if you want to add LaTeX within your code
  extendedchars=true,              % lets you use non-ASCII characters; for 8-bits encodings only, does not work with UTF-8
  frame=single,                    % adds a frame around the code
  keepspaces=true,                 % keeps spaces in text, useful for keeping indentation of code (possibly needs columns=flexible)
  keywordstyle=\color{blue},       % keyword style
  language=Make,                 % the language of the code
  morekeywords={*,...,.include},            % if you want to add more keywords to the set
  numbers=left,                    % where to put the line-numbers; possible values are (none, left, right)
  numbersep=5pt,                   % how far the line-numbers are from the code
  numberstyle=\tiny\color{mygray}, % the style that is used for the line-numbers
  rulecolor=\color{black},         % if not set, the frame-color may be changed on line-breaks within not-black text (e.g. comments (green here))
  showspaces=false,                % show spaces everywhere adding particular underscores; it overrides 'showstringspaces'
  showstringspaces=false,          % underline spaces within strings only
  showtabs=false,                  % show tabs within strings adding particular underscores
  stepnumber=1,                    % the step between two line-numbers. If it's 1, each line will be numbered
  stringstyle=\color{mymauve},     % string literal style
  tabsize=8,                       % sets default tabsize to 2 spaces
  title=\lstname                   % show the filename of files included with \lstinputlisting; also try caption instead of title
}

\addtobeamertemplate{frametitle}{}{%
\begin{tikzpicture}[remember picture,overlay]
	\node[anchor=north east,yshift=-5pt] at (current page.north east) {\includegraphics[scale=0.13]{../ganditemplate/Gandi_logo_black.eps}};
\end{tikzpicture}}

%%\lstset{escapechar=@,style=custommake}

\usepackage{color}
\usepackage{listings}
\usepackage{tipa}
\usepackage{svg}
\definecolor{indigo}{rgb}{0.29, 0.0, 0.51}
\newcommand{\prestitle}{EuroBSDCon 2015}
\lstset{
	numbers=none,
	frame=l,
	aboveskip=0pt, belowskip=0pt
}

\title{Poudrière}
\subtitle{Efficient package building}
\author{Baptiste Daroussin \\ \url{bapt@FreeBSD.org}}
\institute{\includegraphics[scale=0.13]{../ganditemplate/Gandi_logo_black.eps}\hspace{80pt}\includegraphics[scale=0.5]{../ganditemplate/logo.pdf}}
\date{EuroBSDCon 2015 \\ Stockholm \\ Octobre 4th, 2015}

\begin{document}
\begin{frame}[plain]
	\titlepage
\end{frame}

\begin{frame}
	\frametitle{poudri...GNI?: [pu.d\textipa{\textinvscr}i.je\textipa{\textinvscr}]}
	\begin{itemize}
		\item Package building system
		\item Port tester
		\item Quality insurance on packages
		\item Package repository generator
		\item System stress tool
	\end{itemize}
\end{frame}

\begin{frame}
	\frametitle{History}
	\begin{itemize}
		\item 2010-07: Initial work
		\item 2011: Start to be known and used in the french community
		\item 2012-01-31: 1.0 - enter the ports tree
		\item 2012-04-08: 1.2 - limit network on fetch phase
		\item 2012-05-15: 1.3 - pbi support, attract interest of bdrewery@
		\item 2012-08-28: 2.0 - parallel build, ugly html UI (bapt as a designer)
		\item 2012-10-15: 2.2 - Removal of pbi support, support for "sets"
		\item 2013-05-20: 3.0 - ZFS optional, full tmpfs support, nice and reactive web UI (bdrewery designer)
		\item 2013-07: Used in the FreeBSD cluster
		\item 2013-09-22: 3.0.3 - support staging, initial qemu support
		\item 2014-12-04: 3.1.0 - Yet a better web UI
	\end{itemize}
\end{frame}

\begin{frame}
	\frametitle{Design}
	\begin{itemize}
		\item Simple
			\begin{itemize}
				\item Easy to setup:
					\begin{itemize}
						\item only depend on base (by default)
						\item one simple configuration file
						\item few command to prepare the resources
					\end{itemize}
				\item Easy to use
					\begin{itemize}
						\item One single command
						\item Simple subcommands
					\end{itemize}
			\end{itemize}
		\item Resource efficient
			\begin{itemize}
				\item parallel build: by default 1 core == 1 package building
				\item low overhead (resources should be dedicated to build sources not for poudriere itself)
			\end{itemize}
		\item Safe and contained
			\begin{itemize}
				\item all builds in clean jail(8)
				\item only access network during fetch phase
				\item build as regular user
			\end{itemize}
	\end{itemize}
\end{frame}

\begin{frame}
	\frametitle{Design}
	Subcommands:
	\begin{itemize}
		\item bulk: Generate packages for given ports
		\item jail: Manage the jails used by poudriere
		\item ports: Create, update or delete the portstrees
	\end{itemize}
\end{frame}

\begin{frame}[fragile]
	\frametitle{Poudrière: jails}
	\begin{itemize}
			\pause
		\item Fetch release/snapshot/old releases sets
			\pause
		\item Build from sources: git, svn, file, support for branches
			\pause
		\item Full support for src.conf
			\pause
		\item Support for multiple arches (via qemu user emulation)
			\pause
		\item Can have kernel
			\pause
		\item Updateable (via sources or freebsd-update)
	\end{itemize}
	\pause
	Creating a jail
	\begin{lstlisting}
poudriere jail -c -j 102 -v 10.2-RELEASE
\end{lstlisting}
	\pause
	Updating a jail
	\begin{lstlisting}
poudriere jail -u -j 102
\end{lstlisting}
\end{frame}

\begin{frame}[fragile]
	\frametitle{Poudrière: ports}
	\begin{itemize}
			\pause
		\item Fetch from portsnap, git, svn
			\pause
		\item Notion of "default" ports tree
			\pause
	\end{itemize}
	\pause
	Creating a ports tree
	\begin{lstlisting}
poudriere ports -c -p portstree
\end{lstlisting}
	\pause
	Updating a ports tree
	\begin{lstlisting}
poudriere ports -u -p portstree
\end{lstlisting}
\end{frame}

\begin{frame}
	\frametitle{Poudrière: bulk}
	\begin{itemize}
			\pause
		\item Associate a ports tree, a jail and a list of packages to build
			\pause
		\item Massively parallelized (1 port per core, can be find tuned)
			\pause
		\item Support ccache
			\pause
		\item Tuneable via: make.conf (fine grained [<jailname>-[<setname>-[<portstree>-]]]make.conf)
			\pause
		\item Nice WebUI (static files made dynamic via js)
			\pause
		\item Nice cli (with colors and siginfo support)
			\pause
		\item Hooks support
			\pause
		\item Repository generation support (including signature)
			\pause
		\item Default ports tree support
			\pause
		\item Incremental support (aggressive)
			\pause
		\item Restricted support
			\pause
		\item Saving workdir after failure
			\pause
		\item Autodetection of rebuild
			\pause
	\end{itemize}
\end{frame}

\begin{frame}[fragile]
	\frametitle{Poudrière: bulk}
	Building packages:
	\begin{lstlisting}
poudriere bulk -j 102 -f listofpackages.txt
\end{lstlisting}
	Building packages with Q/A:
	\begin{lstlisting}
poudriere bulk -j 102 -t -f listofpackages.txt
\end{lstlisting}
	Building all ports
	\begin{lstlisting}
poudriere bulk -j 102 -a
\end{lstlisting}
	Building all ports with a special "set"
	\begin{lstlisting}
poudriere bulk -z test1 -j 102 -a
\end{lstlisting}
\end{frame}

\begin{frame}
	\frametitle{Poudrière: a stress tool}
	In FreeBSD:
	\begin{itemize}
		\item ZFS deadlocks
		\item tmpfs deadlocks
		\item nullfs deadlocks
		\item tons of fixes in sh(1) in particular regarding job control
		\item highlight contentions
	\end{itemize}
	In Dragonfly:
	\begin{itemize}
		\item Used as a benchmark tool in 2013
		\item Lots of performance improvement between December 26, 2012 and March 15, 2013 (released in 3.4)
		\item Lots of panics fixed
	\end{itemize}
\end{frame}

\begin{frame}
	\frametitle{Poudrière: under the hood}
	\begin{itemize}
			\pause
		\item Mostly coded in sh(1) (clean and maintainable shell is possible!)
			\pause
		\item Small bits in C for perfomances
			\pause
		\item Lots of care made on efficiency:
			\begin{itemize}
				\item avoid subshells as much as possible
				\item parallelize as many things as possible
				\item reuse resources as much as possible
			\end{itemize}
		\item Use filesystem as a Key/Value DB (on tmpfs for speed)
	\end{itemize}
\end{frame}

\begin{frame}
	\frametitle{Poudrière: image (soon)}
	\begin{itemize}
		\item Associate jails, packages and overlays
		\item Able to generates usable images:
			\begin{itemize}
				\item Isos: with or without mfsroot
				\item Usb disk: with or without mfsroot
				\item GPT base firmwares (NanoBSD-like)
				\item plain mfsroot
				\item rawdisk (VMs)
			\end{itemize}
	\end{itemize}
\end{frame}

\begin{frame}
	\frametitle{Poudrière: links}
	\begin{itemize}
		\item https://github.com/freebsd/poudriere/
		\item https://github.com/freebsd/poudriere/wiki
		\item https://www.freebsd.org/doc/handbook/ports-poudriere.html
		\item https://www.freebsd.org/doc/en/books/porters-handbook/testing-poudriere.html
	\end{itemize}
\end{frame}

\begin{frame}
	\frametitle{Questions?}
	\center
	\includegraphics{poudriere.eps}
	\\
	\Huge Thanks
\end{frame}

\end{document}
